\documentclass[preview]{standalone}

\usepackage{amsmath}
\usepackage{amssymb}
\usepackage{stellar}
\usepackage{definitions}
\usepackage{bettelini}

\begin{document}

\id{probability-continuous-random-variables}
\genpage

\section{Absolutely continuous random variables}

\begin{snippetdefinition}{absolutely-continuous-random-variable-definition}{Absolutely continuous random variable}
    Let \((\Omega,\mathcal F,\mathbb P)\) be a \probabilityspace, and let
    \(X\colon\Omega\fromto\realnumbers\) be a real-valued \randomvariable.
    The random variable \(X\) is \emph{absolutely continuous} if its
    \lawofrandomvariable \(\mu_X\) is an \absolutelycontinuousmeasure on
    \((\realnumbers,\mathcal B(\realnumbers))\). A \emph{density} of \(X\)
    is a \probabilitydensity \(f_X\) of \(\mu_X\), so that
    \[
        \mathbb P(X\in A)=\mu_X(A)=\int_A f_X(x)\,dx
    \]
    for every \(A\in\mathcal B(\realnumbers)\).
\end{snippetdefinition}

\begin{snippetproposition}{ac-random-variable-cdf-density-formula}{Distribution function of an absolutely continuous random variable}
    Let \((\Omega,\mathcal F,\mathbb P)\) be a \probabilityspace, and let
    \(X\colon\Omega\fromto\realnumbers\) be an
    \absolutelycontinuousrandomvariable with density \(f_X\). If \(F_X\) is
    the \distributionfunction of \(X\), then
    \[
        F_X(x)=\int_{-\infty}^x f_X(t)\,dt
    \]
    for every \(x\in\realnumbers\). In particular, \(F_X\) is continuous and
    \(f_X=F_X'\) almost everywhere whenever \(F_X\) is differentiable almost
    everywhere and the fundamental theorem for the Lebesgue integral applies.
\end{snippetproposition}

\begin{snippetproof}{ac-random-variable-cdf-density-formula-proof}{ac-random-variable-cdf-density-formula}{Distribution function of an absolutely continuous random variable}
    By the definition of density,
    \[
        F_X(x)=\mathbb P(X\leq x)
        =
        \mu_X((-\infty,x])
        =
        \int_{-\infty}^x f_X(t)\,dt.
    \]
    The continuity follows from absolute continuity of the Lebesgue integral.
    The last statement is the usual almost everywhere differentiation theorem
    for indefinite Lebesgue integrals.
\end{snippetproof}

\begin{snippetproposition}{density-determines-law-not-random-variable}{A density determines the law, not the random variable}
    Let \(f\colon\realnumbers\fromto[0,+\infty]\) be a Borel \function such
    that \(\int_{\realnumbers} f(x)\,dx=1\). Then \(f\) determines a unique
    \probabilitymeasure \(\mu\) on
    \((\realnumbers,\mathcal B(\realnumbers))\), namely
    \[
        \mu(A)=\int_A f(x)\,dx.
    \]
    If \(X\) and \(Y\) are \randomvariable[random variables], possibly
    defined on different \probabilityspace[probability spaces], and both have
    density \(f\), then \(X\eqdist Y\). This does not imply that \(X\) and
    \(Y\) are almost surely equal.
\end{snippetproposition}

\begin{snippetproof}{density-determines-law-not-random-variable-proof}{density-determines-law-not-random-variable}{A density determines the law, not the random variable}
    The first statement is the probability measure induced by a density. If
    \(X\) and \(Y\) both have density \(f\), then their laws agree with
    \(\mu\), hence \(X\eqdist Y\). Almost sure equality is a statement on a
    common probability space; equality in law only compares the induced
    probability measures on \(\realnumbers\).
\end{snippetproof}

\section{Expected value}

\begin{snippetdefinition}{expected-value-random-variable-definition}{Expected value of a random variable}
    Let \((\Omega,\mathcal F,\mathbb P)\) be a \probabilityspace, and let
    \(X\colon\Omega\fromto\realnumbers\) be a real-valued \randomvariable.
    The \emph{expected value} of \(X\) is the Lebesgue integral
    \[
        \expectation[X]=\int_\Omega X(\omega)\,d\mathbb P(\omega),
    \]
    whenever this integral is well-defined in \([-\infty,+\infty]\). The
    random variable \(X\) has a \emph{finite expected value} if
    \[
        \expectation[|X|]<+\infty.
    \]
\end{snippetdefinition}

\begin{snippettheorem}{expected-value-law-formula}{Expected value through the law}
    Let \((\Omega,\mathcal F,\mathbb P)\) be a \probabilityspace, and let
    \(X\colon\Omega\fromto\realnumbers\) be a real-valued \randomvariable
    with \lawofrandomvariable \(\mu_X\). Then \(X\) has a finite
    \expectedvalue if and only if
    \[
        \int_{\realnumbers} |x|\,d\mu_X(x)<+\infty.
    \]
    In this case,
    \[
        \expectation[X]=\int_{\realnumbers} x\,d\mu_X(x).
    \]
\end{snippettheorem}

\begin{snippetproof}{expected-value-law-formula-proof}{expected-value-law-formula}{Expected value through the law}
    This is the change-of-variable formula for the pushforward measure
    \(\mu_X=\mathbb P\circ X^{-1}\). Applying it to \(x\mapsto x\) and to
    \(x\mapsto |x|\) gives both the identity and the finiteness criterion.
\end{snippetproof}

\begin{snippetcorollary}{expected-value-absolutely-continuous-density-formula}{Expected value through a density}
    Let \((\Omega,\mathcal F,\mathbb P)\) be a \probabilityspace, and let
    \(X\colon\Omega\fromto\realnumbers\) be an
    \absolutelycontinuousrandomvariable with density \(f_X\). Then \(X\) has
    a finite \expectedvalue if and only if
    \[
        \int_\realnumbers |x|f_X(x)\,dx<+\infty.
    \]
    In this case,
    \[
        \expectation[X]=\int_\realnumbers xf_X(x)\,dx.
    \]
\end{snippetcorollary}

\begin{snippetproof}{expected-value-absolutely-continuous-density-formula-proof}{expected-value-absolutely-continuous-density-formula}{Expected value through a density}
    Since \(d\mu_X(x)=f_X(x)\,dx\), the formula follows from the expression of
    the \expectedvalue through the \lawofrandomvariable.
\end{snippetproof}

\begin{snippettheorem}{random-variable-change-of-variable-expectation}{Change-of-variable formula for expected value}
    Let \((\Omega,\mathcal F,\mathbb P)\) be a \probabilityspace, let
    \(X\colon\Omega\fromto\realnumbers^n\) be a \randomvariable, and let
    \(g\colon\realnumbers^n\fromto\realnumbers\) be a Borel \function. If
    \(\mu_X\) is the \lawofrandomvariable of \(X\), then \(g(X)\) has a finite
    \expectedvalue if and only if
    \[
        \int_{\realnumbers^n}|g(x)|\,d\mu_X(x)<+\infty.
    \]
    In this case,
    \[
        \expectation[g(X)]
        =
        \int_{\realnumbers^n}g(x)\,d\mu_X(x).
    \]
    If \(\mu_X\) has density \(f_X\) with respect to \(n\)-dimensional
    \lebesguemeasure, then this becomes
    \[
        \expectation[g(X)]
        =
        \int_{\realnumbers^n}g(x)f_X(x)\,dx.
    \]
\end{snippettheorem}

\begin{snippetproof}{random-variable-change-of-variable-expectation-proof}{random-variable-change-of-variable-expectation}{Change-of-variable formula for expected value}
    The \lawofrandomvariable of \(g(X)\) is the pushforward of \(\mu_X\) by
    \(g\). Applying the one-dimensional law formula to \(g(X)\) gives the
    integral with respect to \(\mu_X\). If \(\mu_X\) has density \(f_X\), the
    last formula follows from the density representation of \(\mu_X\).
\end{snippetproof}

\section{Moments}

\begin{snippetdefinition}{finite-moment-random-variable-definition}{Finite moment}
    Let \((\Omega,\mathcal F,\mathbb P)\) be a \probabilityspace, let
    \(X\colon\Omega\fromto\realnumbers\) be a real-valued \randomvariable,
    and let \(p>0\). The random variable \(X\) has a \emph{finite moment of
    order \(p\)} if
    \[
        \expectation[|X|^p]<+\infty.
    \]
\end{snippetdefinition}

\begin{snippetproposition}{higher-moment-implies-lower-moment}{Higher moments imply lower moments}
    Let \((\Omega,\mathcal F,\mathbb P)\) be a \probabilityspace, let
    \(X\colon\Omega\fromto\realnumbers\) be a real-valued \randomvariable,
    and let \(0<q\leq p\). If
    \[
        \expectation[|X|^p]<+\infty,
    \]
    then
    \[
        \expectation[|X|^q]<+\infty.
    \]
\end{snippetproposition}

\begin{snippetproof}{higher-moment-implies-lower-moment-proof}{higher-moment-implies-lower-moment}{Higher moments imply lower moments}
    For every \(x\in\realnumbers\),
    \[
        |x|^q\leq 1+|x|^p.
    \]
    Taking expected values gives
    \[
        \expectation[|X|^q]\leq 1+\expectation[|X|^p]<+\infty.
    \]
\end{snippetproof}

\begin{snippetproposition}{variance-absolutely-continuous-density-formula}{Variance through a density}
    Let \((\Omega,\mathcal F,\mathbb P)\) be a \probabilityspace, and let
    \(X\colon\Omega\fromto\realnumbers\) be an
    \absolutelycontinuousrandomvariable with density \(f_X\). If
    \[
        \int_\realnumbers x^2f_X(x)\,dx<+\infty,
    \]
    then \(X\) has finite variance and
    \[
        \variance(X)
        =
        \int_\realnumbers\left(x-\expectation[X]\right)^2f_X(x)\,dx
        =
        \int_\realnumbers x^2f_X(x)\,dx-\expectation[X]^2.
    \]
\end{snippetproposition}

\begin{snippetproof}{variance-absolutely-continuous-density-formula-proof}{variance-absolutely-continuous-density-formula}{Variance through a density}
    The second moment assumption implies that the first moment is finite by
    the higher-to-lower moment result. Apply the change-of-variable formula to
    \(g(x)=(x-\expectation[X])^2\) and expand the square.
\end{snippetproof}

\end{document}
